\documentclass[10pt,a4paper]{amsart}
\usepackage{amsmath,amssymb,amsthm,graphicx}
% Retained preamble for the Stacks Project chapter snapshot in this directory.
% The chapter's own preamble.tex is a copy of the shared Stacks Project preamble; the block
% below is sliced MECHANICALLY, line for line, from it. Nothing was added, removed or
% redefined; the wrapper class, the xr-hyper/hyperref wiring and the externaldocument list
% are omitted because the assembler supplies the document class.
%
% lines taken from preamble.tex: 3, 5, 9, 11, 12, 156, 157, 158, ... (51 lines total)
%
% package coverage restored from the chapter's own preamble (2026-10-13): xy, xy-2cell

\usepackage[all]{xy}
\xyoption{2cell}
\UseAllTwocells

\newenvironment{reference}{\comment}{\endcomment}
\newenvironment{slogan}{\comment}{\endcomment}
\newenvironment{history}{\comment}{\endcomment}
\theoremstyle{plain}
\newtheorem{theorem}[subsection]{Theorem}
\newtheorem{proposition}[subsection]{Proposition}
\newtheorem{lemma}[subsection]{Lemma}
\theoremstyle{definition}
\newtheorem{definition}[subsection]{Definition}
\newtheorem{example}[subsection]{Example}
\newtheorem{exercise}[subsection]{Exercise}
\newtheorem{situation}[subsection]{Situation}
\theoremstyle{remark}
\newtheorem{remark}[subsection]{Remark}
\newtheorem{remarks}[subsection]{Remarks}
\numberwithin{equation}{subsection}
\def\lim{\mathop{\mathrm{lim}}\nolimits}
\def\colim{\mathop{\mathrm{colim}}\nolimits}
\def\Spec{\mathop{\mathrm{Spec}}}
\def\Hom{\mathop{\mathrm{Hom}}\nolimits}
\def\Ext{\mathop{\mathrm{Ext}}\nolimits}
\def\SheafHom{\mathop{\mathcal{H}\!\mathit{om}}\nolimits}
\def\SheafExt{\mathop{\mathcal{E}\!\mathit{xt}}\nolimits}
\def\Sch{\mathit{Sch}}
\def\Mor{\mathop{\mathrm{Mor}}\nolimits}
\def\Ob{\mathop{\mathrm{Ob}}\nolimits}
\def\Sh{\mathop{\mathit{Sh}}\nolimits}
\def\NL{\mathop{N\!L}\nolimits}
\def\CH{\mathop{\mathrm{CH}}\nolimits}
\def\proetale{{pro\text{-}\acute{e}tale}}
\def\etale{{\acute{e}tale}}
\def\QCoh{\mathit{QCoh}}
\def\Ker{\mathop{\mathrm{Ker}}}
\def\Im{\mathop{\mathrm{Im}}}
\def\Coker{\mathop{\mathrm{Coker}}}
\def\Coim{\mathop{\mathrm{Coim}}}
\DeclareMathSymbol{\boxtimes}{\mathbin}{AMSa}{"02}
\def\QCohstack{\mathcal{QC}\!\mathit{oh}}
\def\Cohstack{\mathcal{C}\!\mathit{oh}}
\def\Spacesstack{\mathcal{S}\!\mathit{paces}}
\def\Quotfunctor{\mathrm{Quot}}
\def\Hilbfunctor{\mathrm{Hilb}}
\def\Curvesstack{\mathcal{C}\!\mathit{urves}}
\def\Polarizedstack{\mathcal{P}\!\mathit{olarized}}
\def\Complexesstack{\mathcal{C}\!\mathit{omplexes}}
\def\Pic{\mathop{\mathrm{Pic}}\nolimits}
\def\Picardstack{\mathcal{P}\!\mathit{ic}}
\def\Picardfunctor{\mathrm{Pic}}
\def\Deformationcategory{\mathcal{D}\!\mathit{ef}}
\title{An essentially finite type local homomorphism is a colimit of essentially finite type ones with quotient transitions}
\author{The Stacks Project Authors}
\date{stacks-project:algebra.tex:local-snapshot-20261012; GNU FDL 1.2 or later}
\begin{document}
\maketitle
\noindent\textit{Editorial note.} All delivered mathematical text is the author's own, copied from the pinned Stacks Project chapter file algebra.tex; no statement or proof step was rewritten, reordered or removed.\par\medskip
\noindent\textit{Editorial note.} Preamble provenance, as required for this source: the frozen snapshot stages the chapter's own preamble file beside it (retained-preamble.tex, 63 lines, 8116 bytes, the same shared Stacks-wide preamble the published chow and etale-cohomology cards use); the first line of the pinned chapter file is \textbackslash{}input\{preamble\}, so no preamble text lives inside the chapter itself. The card takes that file's lines 8-18 (the verbatim package and the three \textbackslash{}comment-based environments reference/slogan/history at lines 9-12, section 13.2, together with the Xy-pic setup) and its lines 154 to the runtime end of the file (the native theorem declarations and mathematical macros). Not one character was changed, no package was added, no macro was defined, and the excerpt serves only as preamble, never as a span coordinate.\par\medskip
\noindent\textit{Editorial note.} The same-chapter pointers are located at their \textbackslash{}label lines inside the pinned chapter file; the cross-chapter pointers are located in the project's own tag list staged in the snapshot (sources/stacks-chow/tags.txt), at the verbatim lines carrying their canonical tags.\par\medskip
\noindent\textit{Editorial note.} The counted spans contain no \textbackslash{}ref and no \textbackslash{}cite, so the card carries no navigation stub and no bibliography.\par\medskip
\noindent\textit{Allowed standard prerequisites (editorial).}\begin{itemize}
\item Essentially of finite type as a localisation of a finitely generated algebra. Used at the start of the counted proof (lines 32514-32516). Supplies the finitely many elements whose localisation is the target.
\item Descending finite data along a colimit presentation of the base. Used in the middle of the counted proof. Moves the finitely many coefficients into one stage of the colimit.
\end{itemize}\medskip
\begin{lemma}
\label{lemma-limit-essentially-finite-type}
Suppose $R \to S$ is a local homomorphism of local rings.
Assume that $S$ is essentially of finite type over $R$.
Then there exists a directed set $(\Lambda, \leq)$, and
a system of local homomorphisms $R_\lambda \to S_\lambda$
of local rings such that
\begin{enumerate}
\item The colimit of the system $R_\lambda \to S_\lambda$
is equal to $R \to S$.
\item Each $R_\lambda$ is essentially of finite type
over $\mathbf{Z}$.
\item Each $S_\lambda$ is essentially of finite type
over $R_\lambda$.
\item For each $\lambda \leq \mu$ the map
$S_\lambda \otimes_{R_\lambda} R_\mu \to S_\mu$
presents $S_\mu$ as the localization of a quotient
of $S_\lambda \otimes_{R_\lambda} R_\mu$.
\end{enumerate}
\end{lemma}

\begin{proof}
Denote $\varphi : R  \to S$ the ring map.
Let $\mathfrak m \subset R$ be the maximal ideal
of $R$ and let $\mathfrak n \subset S$ be the maximal
ideal of $S$. Let $x_1, \ldots, x_n \in S$ be elements such that
$S$ is a localization of the sub $R$-algebra of $S$
generated by $x_1, \ldots, x_n$. In other words, $S$
is a quotient of a localization of the polynomial ring
$R[x_1, \ldots, x_n]$.

\medskip\noindent
Let $\Lambda = \{ A \subset R \mid \# A < \infty\}$
be the set of finite subsets of $R$.
As partial ordering we take the inclusion relation. For each
$\lambda = A \in \Lambda$ we let $R'_\lambda$ be
the sub $\mathbf{Z}$-algebra generated by
$a \in A$, and we let $S'_\lambda$ be the sub
$\mathbf{Z}$-algebra generated by $\varphi(a)$, $a \in A$
and the elements $x_1, \ldots, x_n$. Let $R_\lambda$ be
the localization of $R'_\lambda$ at the prime ideal
$R'_\lambda \cap \mathfrak m$ and let
$S_\lambda$ be the localization of $S'_\lambda$ at
the prime ideal $S'_\lambda \cap \mathfrak n$.
In a picture
$$
\xymatrix{
\varphi(A) \amalg \{x_i\} \ar[r] &
S'_\lambda \ar[r] &
S_\lambda \ar[r] &
S \\
A \ar[r] \ar[u] &
R'_\lambda \ar[r] \ar[u] &
R_\lambda \ar[r] \ar[u] &
R \ar[u]
}
$$
It is clear that if $A \subset B$ corresponds to
$\lambda \leq \mu$ in $\Lambda$, then there are
canonical maps $R_\lambda \to R_\mu$, and $S_\lambda \to S_\mu$
and we obtain a system over the directed set $\Lambda$.

\medskip\noindent
The assertion that $R = \colim R_\lambda$ is clear
because all the maps $R_\lambda \to R$ are injective and
any element of $R$ eventually is in the image. The same
argument works for $S = \colim S_\lambda$.
Assertions (2), (3) are true by construction.
The final assertion holds because clearly
the maps $S'_\lambda \otimes_{R'_\lambda} R'_\mu
\to S'_\mu$ are surjective.
\end{proof}

\end{document}
